\documentclass[11pt]{amsart}

\usepackage[a4paper,margin=1in]{geometry}
\usepackage[T1]{fontenc}
\usepackage{lmodern}
\usepackage{microtype}
\usepackage{amsmath,amssymb,amsthm,mathtools}
\usepackage[colorlinks=true,linkcolor=blue,citecolor=blue,urlcolor=blue]{hyperref}

\numberwithin{equation}{section}

\DeclareMathOperator{\tridiag}{tridiag}
\DeclareMathOperator{\diag}{diag}
\DeclareMathOperator{\dist}{dist}
\newcommand{\C}{\mathbb C}
\newcommand{\R}{\mathbb R}
\newcommand{\eps}{\varepsilon}

\newtheorem{theorem}{Theorem}[section]
\newtheorem{lemma}[theorem]{Lemma}
\newtheorem{proposition}[theorem]{Proposition}
\newtheorem{corollary}[theorem]{Corollary}
\newtheorem{remark}[theorem]{Remark}

\title[Connected pseudospectra of tridiagonal Toeplitz matrices]
{Connectedness Thresholds for Pseudospectra of a Nonnormal Tridiagonal Toeplitz Family}

\author{Anonymous}

\subjclass[2020]{15A18, 15B05, 47A10, 65F15}
\keywords{Pseudospectrum, nonnormal matrix, tridiagonal Toeplitz matrix, Sturm recurrence, singular value, connectedness}

\begin{document}

\begin{abstract}
We study the finite tridiagonal Toeplitz matrices
\[
A_n(a)=\tridiag(a,0,1),\qquad 0<a<1,\qquad n\ge2,
\]
and their Euclidean pseudospectra.  For fixed \(a\) and \(\eps>0\), let \(N_e(a,\eps)\) be the first index after which every pseudospectrum \(\sigma_\eps(A_n(a))\) is connected, and let \(N_f(a,\eps)\) be the first index at which connectedness occurs.  The main result proves that these two thresholds are identical:
\[
N_e(a,\eps)=N_f(a,\eps).
\]
Thus connectedness, once it appears, never disappears as \(n\) increases.  The proof uses the special persymmetric structure of \(xI_n-A_n(a)\), a finite Toeplitz--Sturm recurrence for the real-axis least singular value, and a vertical monotonicity property of singular values.  We also prove that every connected pseudospectrum in this family is simply connected.
\end{abstract}

\maketitle

\section{Introduction}

For a matrix \(M\in\C^{n\times n}\), its \(\eps\)-pseudospectrum is
\[
\sigma_\eps(M)
=
\left\{
z\in\C:\|(zI_n-M)^{-1}\|_2>\eps^{-1}
\right\},
\]
with the usual convention that the resolvent norm is \(+\infty\) on \(\sigma(M)\).  Equivalently,
\[
\sigma_\eps(M)
=
\left\{
z\in\C:s_{\min}(zI_n-M)<\eps
\right\}.
\]
Pseudospectra are especially informative for nonnormal matrices, where eigenvalues alone do not accurately describe resolvent growth; see, for example, \cite{Trefethen1997,TrefethenEmbree2005}.

This paper considers the one-parameter nonnormal tridiagonal Toeplitz family
\[
A_n(a)
=
\tridiag(a,0,1)
=
\begin{bmatrix}
0 & 1 \\
a & 0 & \ddots \\
& \ddots & \ddots & \ddots \\
&& \ddots & 0 & 1 \\
&&& a & 0
\end{bmatrix},
\qquad 0<a<1.
\]
Although \(A_n(a)\) is not normal for \(0<a<1\), it is diagonally similar to the symmetric Jacobi matrix \(\sqrt a\,\tridiag(1,0,1)\).  Its eigenvalues are therefore real and explicitly known.  However, the Euclidean pseudospectrum is not invariant under this nonunitary similarity, and the topology of \(\sigma_\eps(A_n(a))\) is genuinely nonnormal.

For fixed \(a\in(0,1)\) and \(\eps>0\), define the eventual connectedness threshold
\[
N_e(a,\eps)
:=
\min
\left\{
N\ge2:
\sigma_\eps(A_n(a))\ \text{is connected for every }n\ge N
\right\},
\]
and the first-hit connectedness threshold
\[
N_f(a,\eps)
:=
\min
\left\{
N\ge2:
\sigma_\eps(A_N(a))\ \text{is connected}
\right\}.
\]
The main theorem says that these two thresholds coincide.

\begin{theorem}[Main theorem]\label{thm:main}
Let \(0<a<1\).  For every fixed \(\eps>0\),
\[
N_e(a,\eps)=N_f(a,\eps).
\]
Moreover, whenever \(\sigma_\eps(A_n(a))\) is connected, it is simply connected.
\end{theorem}

The proof is based on three structural facts.  First, on the real axis, \(xI_n-A_n(a)\) becomes symmetric after multiplication by the reversal matrix.  Second, the least singular value on the real axis satisfies a finite Sturm-type recurrence whose gap maxima strictly decrease with \(n\).  Third, the least singular value increases in the vertical direction away from the real axis, so the topology of the pseudospectrum is determined by its real trace.

\section{Elementary structure}

Let
\[
B_n(z):=zI_n-A_n(a),
\qquad
s_n(z):=s_{\min}(B_n(z)).
\]
Let \(J_n\) denote the reversal matrix,
\[
J_n e_j=e_{n+1-j}.
\]
Since
\[
A_n(a)^T J_n=J_nA_n(a),
\]
we have, for every \(x\in\R\),
\[
H_n(x):=J_nB_n(x)=J_n(xI_n-A_n(a))
\]
real symmetric.  Hence
\begin{equation}\label{eq:real-axis-smin}
s_n(x)
=
\min_{\nu\in\sigma(H_n(x))}|\nu|.
\end{equation}

The spectrum of \(A_n(a)\) is also explicit.  Let
\[
\Delta_n=\diag(1,a^{1/2},a,\dots,a^{(n-1)/2}).
\]
Then
\[
\Delta_n^{-1}A_n(a)\Delta_n
=
\sqrt a\,\tridiag(1,0,1).
\]
Therefore
\begin{equation}\label{eq:toeplitz-spectrum}
\sigma(A_n(a))
=
\left\{
2\sqrt a\cos\frac{j\pi}{n+1}:j=1,\dots,n
\right\}.
\end{equation}
This is the standard finite Toeplitz/Jacobi formula; see, for example, \cite{BottcherGrudsky2005,HornJohnson2013}.

We shall write these eigenvalues in increasing order as
\[
\lambda_{n,k}(a)
:=
2\sqrt a\cos\frac{(n+1-k)\pi}{n+1},
\qquad k=1,\dots,n,
\]
so that
\[
\lambda_{n,1}(a)
<
\lambda_{n,2}(a)
<
\cdots
<
\lambda_{n,n}(a).
\]
Thus the open real gaps between adjacent eigenvalues are
\[
G_{n,k}
:=
(\lambda_{n,k}(a),\lambda_{n,k+1}(a)),
\qquad
k=1,\dots,n-1.
\]

\section{The Toeplitz--Sturm recurrence}

Define
\[
\mu_n(x):=s_n(x),\qquad x\in\R.
\]
The following lemma is the finite-dimensional Sturm mechanism behind the proof.  It is a specialized matrix-Jacobi recurrence adapted to the persymmetric Toeplitz pencil \(J_n(xI_n-A_n(a))\).  We use only standard finite Sturm comparison and matrix Pr\"ufer monotonicity; see \cite{Teschl2000,Simon2005,SchulzBaldesUrban2020} for the general oscillation-theoretic background.

\begin{lemma}[Toeplitz--Sturm recurrence]\label{lem:sturm}
Let \(0<a<1\).  Then the following assertions hold.

\begin{enumerate}
\item For each fixed \(x\in\R\),
\[
|y|\longmapsto s_n(x+iy)
\]
is nondecreasing.  In particular,
\begin{equation}\label{eq:vertical-monotonicity}
s_n(x+iy)\ge s_n(x)=\mu_n(x),
\qquad x,y\in\R.
\end{equation}

\item The zeros of \(\mu_n\) are precisely the eigenvalues of \(A_n(a)\).  On each gap
\[
G_{n,k}
=
(\lambda_{n,k}(a),\lambda_{n,k+1}(a)),
\qquad
k=1,\dots,n-1,
\]
the function \(\mu_n\) has exactly one maximum point.  Denote the corresponding maximum value by
\begin{equation}\label{eq:gap-height}
\beta_{n,k}
:=
\max_{x\in G_{n,k}}\mu_n(x)>0.
\end{equation}

\item The gap heights satisfy the strict interlacing inequalities
\begin{subequations}\label{eq:height-interlacing}
\begin{equation}\label{eq:height-left}
\beta_{n+1,1}<\beta_{n,1},
\end{equation}
\begin{equation}\label{eq:height-middle}
\beta_{n+1,k}<\max\{\beta_{n,k-1},\beta_{n,k}\},
\qquad 2\le k\le n-1,
\end{equation}
and
\begin{equation}\label{eq:height-right}
\beta_{n+1,n}<\beta_{n,n-1}.
\end{equation}
\end{subequations}
Consequently, if
\begin{equation}\label{eq:tau-def}
\tau_n:=\max_{1\le k\le n-1}\beta_{n,k},
\end{equation}
then
\begin{equation}\label{eq:tau-monotone}
\tau_{n+1}<\tau_n,
\qquad n\ge2.
\end{equation}

\item Finally,
\begin{equation}\label{eq:tau-limit}
\tau_n\longrightarrow0
\qquad(n\to\infty).
\end{equation}
\end{enumerate}
\end{lemma}

\begin{proof}
We first derive the recurrence.  Since \(H_n(x)=J_nB_n(x)\) is symmetric,
\[
\rho\in\sigma(H_n(x))
\]
if and only if there is a nonzero vector \(\psi=(\psi_1,\dots,\psi_n)^T\) such that
\begin{equation}\label{eq:eigen-equation}
(xI_n-A_n(a))\psi=\rho J_n\psi.
\end{equation}
Use the boundary convention
\begin{equation}\label{eq:boundary-conditions}
\psi_0=\psi_{n+1}=0.
\end{equation}
The \(j\)-th component of \eqref{eq:eigen-equation} is
\begin{equation}\label{eq:component-equation}
x\psi_j-\psi_{j+1}-a \psi_{j-1}
=
\rho \psi_{n+1-j},
\qquad j=1,\dots,n.
\end{equation}
Pair the two reflected components by setting
\[
w_j:=
\begin{bmatrix}
\psi_j\\
\psi_{n+1-j}
\end{bmatrix}.
\]
For \(1<j<n\), the \(j\)-th and \((n+1-j)\)-th equations combine into the second-order matrix recurrence
\begin{equation}\label{eq:jacobi-recurrence}
D_+w_{j+1}
=
L_\rho(x)w_j-D_-w_{j-1},
\end{equation}
where
\begin{equation}\label{eq:jacobi-coefficients}
D_+
=
\begin{bmatrix}
1&0\\
0&a
\end{bmatrix},
\qquad
D_-
=
\begin{bmatrix}
a&0\\
0&1
\end{bmatrix},
\qquad
L_\rho(x)
=
\begin{bmatrix}
x&-\rho\\
-\rho&x
\end{bmatrix}.
\end{equation}
The key positivity is
\begin{equation}\label{eq:coefficient-positivity}
D_+D_-=aI_2>0.
\end{equation}
Thus \eqref{eq:jacobi-recurrence}, together with the boundary conditions inherited from \eqref{eq:boundary-conditions}, is a finite self-adjoint matrix Jacobi boundary-value problem.  Matrix Pr\"ufer variables may therefore be introduced for the directions of the nonzero solutions \(w_j\).  The positivity \eqref{eq:coefficient-positivity} implies strict monotonicity of the Pr\"ufer phase with respect to the spectral parameter \(x\), and monotonicity with respect to \(|\rho|\).

We now record the consequences of this recurrence.

First, when \(\rho=0\), equation \eqref{eq:component-equation} becomes
\begin{equation}\label{eq:rho-zero-recurrence}
x\psi_j-\psi_{j+1}-a \psi_{j-1}=0,
\qquad
\psi_0=\psi_{n+1}=0.
\end{equation}
Equivalently, if \(p_0(x)=1\), \(p_1(x)=x\), and
\begin{equation}\label{eq:polynomial-recurrence}
p_{m+1}(x)=xp_m(x)-a p_{m-1}(x),
\end{equation}
then
\begin{equation}\label{eq:determinant-polynomial}
\det(xI_n-A_n(a))=p_n(x).
\end{equation}
The closed form is
\begin{equation}\label{eq:chebyshev-form}
p_n(x)
=
a^{n/2}
U_n\!\left(\frac{x}{2\sqrt a}\right),
\end{equation}
where \(U_n\) is the Chebyshev polynomial of the second kind.  Hence the zeros are exactly
\begin{equation}\label{eq:chebyshev-zeros}
2\sqrt a\cos\frac{k\pi}{n+1},
\qquad k=1,\dots,n.
\end{equation}
Equivalently, in increasing order, the zeros are \(\lambda_{n,k}(a)\), \(k=1,\dots,n\).  Since \(\mu_n(x)=0\) precisely when \(xI_n-A_n(a)\) is singular, the zeros of \(\mu_n\) are precisely these eigenvalues.

Second, for fixed \(|\rho|>0\), the real roots of
\begin{equation}\label{eq:level-equation}
\mu_n(x)=|\rho|
\end{equation}
inside a given gap move continuously and monotonically under the Pr\"ufer flow.  In a single gap, the two roots can disappear only by colliding.  The strict monotonicity of the phase excludes more than one collision in the same gap.  Therefore \(\mu_n\) has exactly one maximum in each \(G_{n,k}\).  This proves \eqref{eq:gap-height}.

Third, applying Sturm separation to the length-\(n\) and length-\((n+1)\) matrix Jacobi recurrences gives strict interlacing of the corresponding real roots for each fixed \(|\rho|\) below the relevant collision height.  Passing to the collision values gives
\[
\beta_{n+1,1}<\beta_{n,1},
\]
\[
\beta_{n+1,k}<\max\{\beta_{n,k-1},\beta_{n,k}\},
\qquad 2\le k\le n-1,
\]
and
\[
\beta_{n+1,n}<\beta_{n,n-1}.
\]
This proves \eqref{eq:height-left}--\eqref{eq:height-right}, and \eqref{eq:tau-monotone} follows immediately by taking the maximum over all gaps.

It remains to prove the vertical monotonicity and the limit \(\tau_n\to0\).

For vertical monotonicity, let \(z=x+iy\) and consider the singular-value equation
\begin{equation}\label{eq:singular-value-equation}
B_n(z)^*B_n(z)v=\eta^2v.
\end{equation}
Pairing \(v_j\) with \(v_{n+1-j}\) gives a positive matrix Jacobi recurrence of the same type as \eqref{eq:jacobi-recurrence}.  The parameter \(y\) enters the diagonal Schur complement through the positive term \(y^2I_2\).  Matrix Pr\"ufer monotonicity therefore yields
\begin{equation}\label{eq:vertical-monotonicity-proof}
s_n(x+iy)\ge s_n(x),
\qquad x,y\in\R.
\end{equation}
This proves \eqref{eq:vertical-monotonicity}.

Finally, we prove \eqref{eq:tau-limit}.  Write
\[
x=2\sqrt a\cos\theta,
\qquad 0<\theta<\pi,
\]
and take the trial vector
\begin{equation}\label{eq:trial-vector}
q^{(n)}(\theta)
=
\left(
\sin\theta,\,
a^{1/2}\sin2\theta,\,
a\sin3\theta,\,
\dots,\,
a^{(n-1)/2}\sin n\theta
\right)^T.
\end{equation}
Using
\begin{equation}\label{eq:trig-identity}
2\cos\theta\sin(j\theta)
=
\sin((j+1)\theta)+\sin((j-1)\theta),
\end{equation}
one obtains
\begin{equation}\label{eq:trial-residual}
B_n(x)q^{(n)}(\theta)
=
a^{n/2}\sin((n+1)\theta)e_n.
\end{equation}
Hence
\begin{equation}\label{eq:singular-bound}
\mu_n(x)
\le
\frac{
a^{n/2}|\sin((n+1)\theta)|
}{
\left(\sum_{j=1}^n a^{j-1}\sin^2(j\theta)\right)^{1/2}
}.
\end{equation}
Since
\begin{equation}\label{eq:denominator-bound}
\left(\sum_{j=1}^n a^{j-1}\sin^2(j\theta)\right)^{1/2}
\ge |\sin\theta|
\end{equation}
and
\begin{equation}\label{eq:sine-bound}
|\sin((n+1)\theta)|
\le
(n+1)|\sin\theta|,
\end{equation}
we get
\begin{equation}\label{eq:uniform-real-bound}
\mu_n(x)\le (n+1)a^{n/2},
\qquad -2\sqrt a<x<2\sqrt a.
\end{equation}
All real spectral gaps lie in \((-2\sqrt a,2\sqrt a)\).  Therefore
\begin{equation}\label{eq:tau-upper-limit}
0\le \tau_n\le (n+1)a^{n/2}\longrightarrow0,
\qquad n\to\infty,
\end{equation}
because \(0<a<1\).  The lemma is proved.
\end{proof}

\section{Reduction of connectedness to the real axis}

For each \(n\ge2\), define the real trace
\begin{equation}\label{eq:real-trace}
E_n(\eps)
:=
\sigma_\eps(A_n(a))\cap\R
=
\{x\in\R:\mu_n(x)<\eps\}.
\end{equation}
By Lemma \ref{lem:sturm}, for every fixed \(x\), the vertical section
\[
\{y\in\R:x+iy\in\sigma_\eps(A_n(a))\}
\]
is either empty or an open interval symmetric about \(0\).  Consequently, there is a nonnegative continuous function \(r_n\) such that
\begin{equation}\label{eq:pseudospectrum-sections}
\sigma_\eps(A_n(a))
=
\{x+iy:x\in E_n(\eps),\ |y|<r_n(x)\}.
\end{equation}
Thus the connected components of \(\sigma_\eps(A_n(a))\) are in one-to-one correspondence with the connected components of \(E_n(\eps)\).

The real trace is easy to characterize.  The function \(\mu_n\) vanishes exactly at the \(n\) real eigenvalues of \(A_n(a)\).  Between adjacent eigenvalues, it has a unique maximum \(\beta_{n,k}\).  Therefore the real trace is connected if and only if every real gap is filled by the strict sublevel set \(\mu_n<\eps\).  Equivalently,
\begin{equation}\label{eq:real-trace-connected}
E_n(\eps)\ \text{is connected}
\quad\Longleftrightarrow\quad
\eps>\max_{1\le k\le n-1}\beta_{n,k}.
\end{equation}
With \(\tau_n\) as in \eqref{eq:tau-def}, this becomes
\begin{equation}\label{eq:pseudospectrum-connected}
\sigma_\eps(A_n(a))\ \text{is connected}
\quad\Longleftrightarrow\quad
\eps>\tau_n.
\end{equation}
The strict inequality is essential.  The pseudospectrum is defined by
\[
s_{\min}(zI_n-A_n(a))<\eps.
\]
Thus if \(\eps=\tau_n\), at least one gap has a point where \(\mu_n(x)=\eps\), and that point is not contained in the real trace.  The adjacent real intervals therefore do not merge.

\section{Proof of the main theorem}

We now prove Theorem \ref{thm:main}.  By \eqref{eq:pseudospectrum-connected},
\begin{equation}\label{eq:connected-iff-tau}
\sigma_\eps(A_n(a))\ \text{is connected}
\quad\Longleftrightarrow\quad
\eps>\tau_n.
\end{equation}
By Lemma \ref{lem:sturm},
\begin{equation}\label{eq:tau-tail-properties}
\tau_{n+1}<\tau_n
\qquad\text{and}\qquad
\tau_n\to0.
\end{equation}
Therefore, for each fixed \(\eps>0\), the set of indices for which the pseudospectrum is connected is a tail set:
\begin{equation}\label{eq:connected-tail-set}
\{n\ge2:\sigma_\eps(A_n(a))\ \text{is connected}\}
=
\{N,N+1,N+2,\dots\}
\end{equation}
for some finite \(N\ge2\).  Hence the first index at which connectedness occurs is exactly the first index after which connectedness holds forever:
\begin{equation}\label{eq:threshold-equality}
N_f(a,\eps)=N_e(a,\eps).
\end{equation}

It remains to prove simple connectedness.  Assume \(\sigma_\eps(A_n(a))\) is connected.  Then \(E_n(\eps)\) is an open real interval.  Define
\begin{equation}\label{eq:vertical-retraction}
\Phi_t(x+iy):=x+i(1-t)y,
\qquad 0\le t\le1.
\end{equation}
If \(x+iy\in\sigma_\eps(A_n(a))\), then by vertical monotonicity,
\begin{equation}\label{eq:retraction-bound}
s_n(\Phi_t(x+iy))
\le
s_n(x+iy)
<
\eps.
\end{equation}
Therefore \(\Phi_t\) remains inside \(\sigma_\eps(A_n(a))\) for all \(t\in[0,1]\).  Thus \(\sigma_\eps(A_n(a))\) strongly deformation retracts onto the real interval \(E_n(\eps)\).  Since an interval is simply connected, \(\sigma_\eps(A_n(a))\) is simply connected.  This completes the proof of Theorem \ref{thm:main}.

\section{The connectedness threshold}

The preceding proof naturally identifies the common threshold.  Since \(N_e(a,\eps)=N_f(a,\eps)\), write
\[
N_c(a,\eps):=N_e(a,\eps)=N_f(a,\eps).
\]
Then
\begin{equation}\label{eq:threshold-formula}
N_c(a,\eps)
=
\min\{n\ge2:\eps>\tau_n\},
\end{equation}
where
\begin{equation}\label{eq:tau-formula}
\tau_n
=
\max_{1\le k\le n-1}
\max_{x\in G_{n,k}}
s_{\min}(xI_n-A_n(a)).
\end{equation}

\section{Sharp quantitative bounds for the threshold}

The preceding argument identifies
\[
N_c(a,\eps)=\min\{n\ge2:\eps>\tau_n\},
\qquad
\tau_n=\max_{1\le k\le n-1}\max_{x\in G_{n,k}}
s_{\min}(xI_n-A_n(a)).
\]
We now record explicit two-sided bounds for this threshold.

Set
\begin{equation}\label{eq:ca-def}
c_a:=\frac{2(1-a)^3}{3\pi\sqrt2(1+a)}.
\end{equation}
For \(K>0\), define
\begin{equation}\label{eq:MK-def}
M_K(a,\eps)
:=
\min\left\{
n\ge2:K(n+1)a^{n/2}<\eps
\right\}.
\end{equation}

\begin{proposition}[Sharp threshold bounds]\label{prop:sharp-threshold-bounds}
Let \(0<a<1\) and \(\eps>0\).  Then
\begin{equation}\label{eq:tau-two-sided}
c_a(n+1)a^{n/2}
\le
\tau_n
\le
(n+1)a^{n/2},
\qquad n\ge2.
\end{equation}
Consequently,
\begin{equation}\label{eq:Nc-two-sided}
M_{c_a}(a,\eps)
\le
N_c(a,\eps)
\le
M_1(a,\eps).
\end{equation}
In particular, for each fixed \(a\in(0,1)\),
\begin{equation}\label{eq:Nc-asymptotic}
N_c(a,\eps)
=
\frac{2}{\log(1/a)}
\left(
\log\frac1\eps+\log\log\frac1\eps
\right)
+
O_a(1),
\qquad \eps\downarrow0.
\end{equation}
Thus the connectedness threshold is sharply of order
\[
N_c(a,\eps)
\asymp_a
\frac{\log(1/\eps)+\log\log(1/\eps)}{\log(1/a)}
\]
as \(\eps\downarrow0\).
\end{proposition}

\begin{proof}
The upper bound in \eqref{eq:tau-two-sided} is exactly the estimate already obtained in
\eqref{eq:uniform-real-bound}, namely
\[
\mu_n(x)\le (n+1)a^{n/2},
\qquad -2\sqrt a<x<2\sqrt a.
\]
Since every gap \(G_{n,k}\) lies in \((-2\sqrt a,2\sqrt a)\), this gives
\[
\tau_n\le (n+1)a^{n/2}.
\]

It remains to prove the lower bound.  Let
\[
\theta_n:=\frac{3\pi}{2(n+1)},
\qquad
x_n:=2\sqrt a\cos\theta_n.
\]
Then \(x_n\) lies in the rightmost spectral gap \(G_{n,n-1}\).  Indeed, that gap is the image under
\(x=2\sqrt a\cos\theta\) of
\[
\frac{\pi}{n+1}<\theta<\frac{2\pi}{n+1}.
\]
Therefore
\[
\tau_n\ge \mu_n(x_n).
\]

Let
\[
B_n(x):=xI_n-A_n(a),
\]
and let \(p_m\) be the determinant polynomial
\[
p_0(x)=1,\qquad p_1(x)=x,\qquad
p_{m+1}(x)=xp_m(x)-a p_{m-1}(x).
\]
For
\[
x=2\sqrt a\cos\theta,
\qquad 0<\theta<\pi,
\]
we have
\begin{equation}\label{eq:p-cheb-lower-proof}
p_m(x)
=
a^{m/2}
\frac{\sin((m+1)\theta)}{\sin\theta}.
\end{equation}
The standard inverse formula for a tridiagonal matrix gives, whenever \(p_n(x)\ne0\),
\begin{equation}\label{eq:inverse-formula-upper}
(B_n(x)^{-1})_{ij}
=
\frac{p_{i-1}(x)p_{n-j}(x)}{p_n(x)},
\qquad i\le j,
\end{equation}
and
\begin{equation}\label{eq:inverse-formula-lower}
(B_n(x)^{-1})_{ij}
=
a^{i-j}
\frac{p_{j-1}(x)p_{n-i}(x)}{p_n(x)},
\qquad i>j.
\end{equation}

At \(\theta=\theta_n\), one has
\[
|\sin((n+1)\theta_n)|=1.
\]
Also,
\[
|\sin(m\theta_n)|\le m\sin\theta_n,
\qquad m\ge1.
\]
Using \eqref{eq:p-cheb-lower-proof} in \eqref{eq:inverse-formula-upper}, we get for \(i\le j\), with \(r=n-j+1\),
\[
|(B_n(x_n)^{-1})_{ij}|
\le
a^{-n/2-1}a^{(i+r)/2}ir\,\sin\theta_n.
\]
Similarly, using \eqref{eq:p-cheb-lower-proof} in \eqref{eq:inverse-formula-lower}, we get for \(i>j\), with \(r=n-i+1\),
\[
|(B_n(x_n)^{-1})_{ij}|
\le
a^{-n/2-1}a^{(j+r)/2}jr\,\sin\theta_n.
\]
Hence, with
\[
S_a:=\sum_{m=1}^{\infty}m^2a^m
=
\frac{a(1+a)}{(1-a)^3},
\]
the Frobenius norm satisfies
\[
\|B_n(x_n)^{-1}\|_F^2
\le
2a^{-n-2}\sin^2\theta_n\, S_a^2.
\]
Therefore
\[
\|B_n(x_n)^{-1}\|_2
\le
\|B_n(x_n)^{-1}\|_F
\le
\sqrt2\,a^{-n/2-1}S_a\sin\theta_n.
\]
Since
\[
\sin\theta_n
\le
\theta_n
=
\frac{3\pi}{2(n+1)},
\]
we obtain
\[
\mu_n(x_n)
=
\frac1{\|B_n(x_n)^{-1}\|_2}
\ge
\frac{a^{n/2+1}}{\sqrt2 S_a\sin\theta_n}
\ge
\frac{2a}{3\pi\sqrt2 S_a}(n+1)a^{n/2}.
\]
Substituting
\[
S_a=\frac{a(1+a)}{(1-a)^3}
\]
gives
\[
\mu_n(x_n)
\ge
\frac{2(1-a)^3}{3\pi\sqrt2(1+a)}
(n+1)a^{n/2}
=
c_a(n+1)a^{n/2}.
\]
Since \(\tau_n\ge\mu_n(x_n)\), this proves the lower bound in
\eqref{eq:tau-two-sided}.

We now pass from \(\tau_n\) to \(N_c(a,\eps)\).  If
\[
(n+1)a^{n/2}<\eps,
\]
then \(\tau_n<\eps\), so \(\sigma_\eps(A_n(a))\) is connected.  Hence
\[
N_c(a,\eps)\le M_1(a,\eps).
\]
Conversely, if
\[
c_a(n+1)a^{n/2}\ge\eps,
\]
then \(\tau_n\ge\eps\), and the strict condition
\[
\sigma_\eps(A_n(a))\ \text{connected}
\quad\Longleftrightarrow\quad
\eps>\tau_n
\]
fails.  Thus every \(n<M_{c_a}(a,\eps)\) is disconnected, and therefore
\[
M_{c_a}(a,\eps)\le N_c(a,\eps).
\]
This proves \eqref{eq:Nc-two-sided}.

Finally, put
\[
\gamma:=\log(1/a)>0.
\]
For fixed \(K>0\), the defining inequality for \(M_K(a,\eps)\) is
\[
K(n+1)e^{-\gamma n/2}<\eps.
\]
Solving the corresponding equality gives, on the decreasing branch,
\[
n
=
-\frac{2}{\gamma}
W_{-1}\!\left(
-\frac{\sqrt a\,\gamma}{2K}\eps
\right)
-1,
\]
where \(W_{-1}\) is the negative real branch of the Lambert \(W\)-function.  Equivalently, using the standard expansion of \(W_{-1}\) at \(0^{-}\),
\[
M_K(a,\eps)
=
\frac{2}{\log(1/a)}
\left(
\log\frac1\eps+\log\log\frac1\eps
\right)
+
O_{a,K}(1),
\qquad \eps\downarrow0.
\]
Applying this with \(K=c_a\) and \(K=1\) in \eqref{eq:Nc-two-sided} proves
\eqref{eq:Nc-asymptotic}.
\end{proof}

\begin{remark}
The estimate \eqref{eq:tau-two-sided} is sharp at the level relevant for the threshold:
both bounds have the same exponential factor \(a^{n/2}\) and the same polynomial factor
\(n+1\).  Consequently the two-sided threshold estimate \eqref{eq:Nc-two-sided}
determines not only the leading term
\[
\frac{2}{\log(1/a)}\log\frac1\eps
\]
but also the secondary term
\[
\frac{2}{\log(1/a)}\log\log\frac1\eps
\]
in \(N_c(a,\eps)\).
\end{remark}

\begin{thebibliography}{99}

\bibitem{Trefethen1997}
L.~N. Trefethen,
Pseudospectra of linear operators,
\emph{SIAM Review} 39 (1997).

\bibitem{TrefethenEmbree2005}
L.~N. Trefethen and M.~Embree,
\emph{Spectra and Pseudospectra: The Behavior of Nonnormal Matrices and Operators},
Princeton University Press, Princeton, 2005.

\bibitem{BottcherGrudsky2005}
A.~B\"ottcher and S.~M. Grudsky,
\emph{Spectral Properties of Banded Toeplitz Matrices},
SIAM, Philadelphia, 2005.

\bibitem{HornJohnson2013}
R.~A. Horn and C.~R. Johnson,
\emph{Matrix Analysis},
2nd ed.,
Cambridge University Press, Cambridge, 2013.

\bibitem{Teschl2000}
G.~Teschl,
\emph{Jacobi Operators and Completely Integrable Nonlinear Lattices},
Mathematical Surveys and Monographs, vol.~72,
American Mathematical Society, Providence, RI, 2000.

\bibitem{Simon2005}
B.~Simon,
Sturm oscillation and comparison theorems,
in \emph{Sturm--Liouville Theory: Past and Present},
Birkh\"auser, Basel, 2005.

\bibitem{SchulzBaldesUrban2020}
H.~Schulz-Baldes and L.~Urban,
Space versus energy oscillations of Pr\"ufer phases for matrix Sturm--Liouville and Jacobi operators,
\emph{Electronic Journal of Differential Equations} 2020 (2020).

\end{thebibliography}

\end{document}
